% !TeX root=main.tex
% در این فایل، عنوان پایان‌نامه، مشخصات خود، متن تقدیمی‌، ستایش، سپاس‌گزاری و چکیده پایان‌نامه را به فارسی، وارد کنید.
% توجه داشته باشید که جدول حاوی مشخصات پروژه/پایان‌نامه/رساله و همچنین، مشخصات داخل آن، به طور خودکار، درج می‌شود.
%%%%%%%%%%%%%%%%%%%%%%%%%%%%%%%%%%%%
% دانشگاه خود را وارد کنید
\university{دانشگاه تهران}
% پردیس دانشگاهی خود را اگر نیاز است وارد کنید (مثال: فنی، علوم پایه، علوم انسانی و ...)
\college{پردیس دانشکده‌های فنی}
% دانشکده، آموزشکده و یا پژوهشکده  خود را وارد کنید
\faculty{دانشکده علوم مهندسی}
% گروه آموزشی خود را وارد کنید (در صورت نیاز)
\department{گروه الگوریتم‌ها و محاسبات}
% رشته تحصیلی خود را وارد کنید
\subject{مهندسی کامپیوتر}
% گرایش خود را وارد کنید
\field{الگوریتم‌ها و محاسبات}
% عنوان پایان‌نامه را وارد کنید
\title{طراحی و بهینه‌سازی الگوریتم اجماع مقیاس‌پذیر برای سیستم‌های بلاکچین}
% نام استاد(ان) راهنما را وارد کنید
\firstsupervisor{دکتر علی کمندی}
\firstsupervisorrank{استاد}
%\secondsupervisor{دکتر سید مرتضی امینی}
%\secondsupervisorrank{استادیار}
% نام استاد(دان) مشاور را وارد کنید. چنانچه استاد مشاور ندارید، دستورات پایین را غیرفعال کنید.
%\firstadvisor{دکتر علی فهیم}
%\firstadvisorrank{استادیار}
%\secondadvisor{دکتر مشاور دوم}
% نام داوران داخلی و خارجی خود را وارد نمایید.
\internaljudge{دکتر علی کمندی}
\internaljudgerank{استاد}
\externaljudge{دکتر محمد رضایی}
\externaljudgerank{دانشیار}
\externaljudgeuniversity{دانشگاه تهران}
% نام نماینده کمیته تحصیلات تکمیلی در دانشکده \ گروه
\graduatedeputy{دکتر علی کمندی}
\graduatedeputyrank{استاد}
% نام دانشجو را وارد کنید
\name{سجاد}
% نام خانوادگی دانشجو را وارد کنید
\surname{خشدونی فراهانی}
% شماره دانشجویی دانشجو را وارد کنید
\studentID{-}
% تاریخ پایان‌نامه را وارد کنید
\thesisdate{آذر 1404}
% به صورت پیش‌فرض برای پایان‌نامه‌های کارشناسی تا دکترا به ترتیب از عبارات «پروژه»، «پایان‌نامه» و «رساله» استفاده می‌شود؛ اگر  نمی‌پسندید هر عنوانی را که مایلید در دستور زیر قرار داده و آنرا از حالت توضیح خارج کنید.
%\projectLabel{پایان‌نامه}

% به صورت پیش‌فرض برای عناوین مقاطع تحصیلی کارشناسی تا دکترا به ترتیب از عبارت «کارشناسی»، «کارشناسی ارشد» و «دکتری» استفاده می‌شود؛ اگر نمی‌پسندید هر عنوانی را که مایلید در دستور زیر قرار داده و آنرا از حالت توضیح خارج کنید.
%\degree{}
%%%%%%%%%%%%%%%%%%%%%%%%%%%%%%%%%%%%%%%%%%%%%%%%%%%%
%% پایان‌نامه خود را تقدیم کنید! %%
\dedication
{
{\Large تقدیم به:}\\
\begin{flushleft}{
	\huge
	\vspace{7mm}
	پدر و مادر عزیزم\\
	و تمام کسانی که به آینده فناوری بلاکچین اعتقاد دارند
}
\end{flushleft}
}
%% متن قدردانی %%
%% ترجیحا با توجه به ذوق و سلیقه خود متن قدردانی را تغییر دهید.
\acknowledgement{
سپاس خداوندگار حکیم را که با لطف بی‌کران خود، آدمی را به زیور عقل آراست و به او توان تفکر و نوآوری عطا فرمود.

در آغاز، وظیفه خود می‌دانم از زحمات بی‌دریغ استاد راهنمای گرامی، جناب آقای دکتر علی کمندی، صمیمانه تشکر و قدردانی کنم که در طول انجام این پایان‌نامه با نهایت صبوری همواره راهنما و مشوق من بودند و قطعاً بدون راهنمایی‌های ارزنده ایشان، این تحقیق به سرانجام نمی‌رسید.

همچنین از جامعه توسعه‌دهندگان متن‌باز و جامعه بلاکچین که با تلاش‌های خود راه را برای چنین تحقیقاتی هموار کرده‌اند، سپاسگزارم.

در پایان، بوسه می‌زنم بر دستان خداوندگاران مهر و مهربانی، پدر و مادر عزیزم که با حمایت‌های بی‌دریغ‌شان، همواره بهترین پشتیبان من بوده‌اند. همچنین از خانواده عزیزم به پاس عاطفه سرشار و گرمای امیدبخش وجودشان، سپاسگزاری می‌کنم.
}
%%%%%%%%%%%%%%%%%%%%%%%%%%%%%%%%%%%%
%چکیده پایان‌نامه را وارد کنید
\fa-abstract{
مقیاس‌پذیری یکی از مهم‌ترین چالش‌های فناوری بلاکچین است که مانع اصلی استفاده گسترده از آن در برنامه‌های کاربردی واقعی به شمار می‌رود. بلاکچین بیتکوین به عنوان پرکاربردترین بلاکچین، تنها قادر به پردازش حدود 7 تراکنش در ثانیه است که برای استفاده در مقیاس جهانی کافی نیست.

در این پایان‌نامه، معماری نوآورانه‌ای برای حل مشکل مقیاس‌پذیری بلاکچین ارائه می‌شود. معماری پیشنهادی بر اساس ساختار زنجیره اصلی {(}\lr{Master Chain}{)} و زنجیره‌های جانبی {(}\lr{Side Chains}{)} است که امکان پردازش موازی تراکنش‌ها را فراهم می‌کند. زنجیره اصلی با استفاده از الگوریتم اثبات کار {(}\lr{PoW}{)} و با زمان تولید بلوک 10 دقیقه، امنیت شبکه را تضمین می‌کند، در حالی که زنجیره‌های جانبی با زمان تولید بلوک 1 دقیقه، تراکنش‌ها را به صورت موازی پردازش می‌کنند.

نتایج آزمایشات نشان می‌دهد که با استفاده از 50 زنجیره جانبی، تعداد تراکنش در ثانیه {(}\lr{TPS}{)} به 85 می‌رسد که معادل 12.14 برابر بهبود نسبت به بیتکوین است. همچنین، با بهینه‌سازی پارامترها، این عدد می‌تواند تا 3400 TPS افزایش یابد که معادل 488 برابر بیتکوین است. معماری پیشنهادی دارای قابلیت مقیاس‌پذیری افقی است، به این معنا که با افزودن زنجیره‌های جانبی جدید، ظرفیت شبکه به صورت خطی افزایش می‌یابد. این ویژگی به همراه حفظ امنیت، غیرمتمرکز بودن و کارایی انرژی بهتر، معماری پیشنهادی را به راه‌حل مناسبی برای توسعه برنامه‌های بلاکچینی در مقیاس بزرگ تبدیل می‌کند.

پیاده‌سازی این معماری در زبان برنامه‌نویسی Rust انجام شده و کد منبع به صورت متن‌باز در دسترس است. نتایج این تحقیق می‌تواند پایه‌ای برای توسعه نسل جدیدی از بلاکچین‌های مقیاس‌پذیر باشد.
}
% کلمات کلیدی پایان‌نامه را وارد کنید
\keywords{بلاکچین، مقیاس‌پذیری، زنجیره اصلی، زنجیره جانبی، اثبات کار، تراکنش در ثانیه، بیتکوین، سیستم‌های توزیع شده}
% انتهای وارد کردن فیلد‌ها
%%%%%%%%%%%%%%%%%%%%%%%%%%%%%%%%%%%%%%%%%%%%%%%%%%%%%%
% ابتدای درج صفحات مختلف
% جلوگیری از ایجاد صفحات خالی در حالت دورو
\let\cleardoublepage\clearpage
\titlePage
% بررسی حالت پیش‌نویس
\ifoptiondraft{}{% 
	\besmPage
	\davaranPage
%%%%%%%%%%%%%%%%%%%%%%%%%%%%
% در این قسمت اسامی اساتید راهنما، مشاور و داور باید به صورت دستی وارد شوند
%\renewcommand{\arraystretch}{1.2}


%%%%%%%%%%%%%%%%%%%%%%%%%%%
	\esalatPage
	\mojavezPage
% چنانچه مایل به چاپ صفحات «تقدیم»، «نیایش» و «سپاس‌گزاری» در خروجی نیستید، خط‌های زیر را با گذاشتن ٪  در ابتدای آنها غیرفعال کنید.
	\taghdimPage
	\ghadrdaniPage
} % end ifoptiondraft
\abstractPage
% شروع درج فهرست‌ها
\clearpage
\pagenumbering{harfi} % آ، ب، ...
\tableofcontents \newpage
% بررسی حالت پیش‌نویس برای بقیه فهرست‌ها
\ifoptiondraft{
	\listoftodos
}{%
	\listoffigures \newpage
	\listoftables  \newpage
	%\addcontentsline{toc}{chapter}{\listalgorithmname}
	%\listofalgorithms %\newpage
	\addcontentsline{toc}{chapter}{\lstlistlistingname}
	\lstlistoflistings \newpage
	%\printacronyms
} % end ifoptiondraft

\pagestyle{fancy}